\documentclass[aspectratio=169]{beamer}

\usepackage{fontspec}
\usepackage[french]{babel}
\usepackage{hyperref}
\usepackage{listings}
\usepackage{xcolor}

\usetheme{Madrid}
\usecolortheme{default}

\definecolor{codebg}{RGB}{245,245,245}
\lstset{
  basicstyle=\ttfamily\footnotesize,
  backgroundcolor=\color{codebg},
  breaklines=true,
  frame=single,
  keywordstyle=\color{blue},
  commentstyle=\color{gray}
}

\title{Quiz de Culture Générale}
\subtitle{Présentation technique du projet}
\author{Antoine Aubry}
\institute{HEIG-VD}
\date{\today}

\begin{document}

% ================================
\begin{frame}
  \titlepage
\end{frame}

% ================================
\begin{frame}{Plan}
  \tableofcontents
\end{frame}

% ================================
\section{Présentation du projet}

\begin{frame}{Le projet en bref}
  \begin{itemize}
    \item Application web de quiz de culture générale (SPA)
    \item L'utilisateur choisit une \textbf{catégorie} et une \textbf{difficulté}
    \item 10 questions à choix multiple, \textbf{temps limité} par question
    \item Score final + \textbf{récapitulatif détaillé} des réponses
    \item Personnalisation : \textbf{thème} clair/sombre, \textbf{langue} FR/EN, support \textbf{manette}
  \end{itemize}
\end{frame}

\begin{frame}[fragile]{Architecture globale}
  \begin{center}
  \begin{lstlisting}
Utilisateur <-> React (SPA) <-> Open Trivia DB (API REST)
                  |
                  +-- localStorage (theme, langue)
  \end{lstlisting}
  \end{center}
  \vspace{0.5cm}
  \begin{itemize}
    \item Application 100\% \textbf{côté client} (pas de backend propre)
    \item Toutes les données viennent d'une API publique externe
    \item Les préférences (thème/langue) sont persistées localement
  \end{itemize}
\end{frame}

% ================================
\section{Frontend : React}

\begin{frame}{React}
  \textbf{Qu'est-ce que c'est ?} Bibliothèque JavaScript pour construire des interfaces utilisateur basées sur des composants.

  \vspace{0.3cm}
  \textbf{Pourquoi ce choix ?}
  \begin{itemize}
    \item Très répandu, large écosystème et documentation
    \item Modèle par composants : découpage naturel (Quiz, QuestionCard, Setup, Home...)
    \item Hooks (\texttt{useState}, \texttt{useEffect}, \texttt{useContext}) pour gérer état et effets de bord
  \end{itemize}
\end{frame}

\begin{frame}{React : avantages / inconvénients}
  \begin{columns}
    \begin{column}{0.5\textwidth}
      \textbf{Avantages}
      \begin{itemize}
        \item Re-rendu efficace via le DOM virtuel
        \item Composants réutilisables et testables
        \item Communauté énorme, beaucoup de ressources
      \end{itemize}
    \end{column}
    \begin{column}{0.5\textwidth}
      \textbf{Inconvénients}
      \begin{itemize}
        \item Courbe d'apprentissage (JSX, hooks, cycle de vie)
        \item Beaucoup de \og boilerplate \fg{} pour des apps simples
        \item Pièges classiques : dépendances de \texttt{useEffect}, \textit{stale state}
      \end{itemize}
    \end{column}
  \end{columns}

  \vspace{0.3cm}
  \textbf{Ce que j'ai appris} : la gestion fine des \texttt{useEffect} (ex. bug du timer où un effet lisait un état \og périmé \fg{} lors du changement de question).
\end{frame}

% ================================
\section{Navigation : React Router DOM}

\begin{frame}{React Router DOM}
  \textbf{Qu'est-ce que c'est ?} Librairie de routage côté client pour les SPA React.

  \vspace{0.3cm}
  \textbf{Pourquoi ce choix ?}
  \begin{itemize}
    \item Standard de facto pour React
    \item Permet 3 pages distinctes (\texttt{/}, \texttt{/setup}, \texttt{/quiz}) sans rechargement
    \item \texttt{useSearchParams} : transmettre catégorie/difficulté via l'URL
    \item \texttt{useNavigate} : redirections programmatiques (rejouer, changer de paramètres)
  \end{itemize}

  \vspace{0.3cm}
  \begin{columns}
    \begin{column}{0.5\textwidth}
      \textbf{Avantages}
      \begin{itemize}
        \item Navigation fluide, état conservé
        \item Paramètres d'URL = configuration partageable
      \end{itemize}
    \end{column}
    \begin{column}{0.5\textwidth}
      \textbf{Inconvénients}
      \begin{itemize}
        \item API qui a beaucoup changé entre versions
        \item Concepts supplémentaires à maîtriser (Routes, Outlet...)
      \end{itemize}
    \end{column}
  \end{columns}
\end{frame}

% ================================
\section{Requêtes HTTP : Axios}

\begin{frame}{Axios}
  \textbf{Qu'est-ce que c'est ?} Client HTTP basé sur les promesses, pour appeler des API REST.

  \vspace{0.3cm}
  \textbf{Pourquoi ce choix ?}
  \begin{itemize}
    \item Syntaxe simple (\texttt{axios.get(url).then(...)})
    \item Gestion d'erreurs centralisée (\texttt{.catch})
    \item Alternative plus pratique que \texttt{fetch} natif (pas de \texttt{.json()} manuel, gestion des erreurs HTTP)
  \end{itemize}

  \vspace{0.3cm}
  \begin{columns}
    \begin{column}{0.5\textwidth}
      \textbf{Avantages}
      \begin{itemize}
        \item Code lisible et concis
        \item Gestion fine des codes d'erreur (ex. 429)
      \end{itemize}
    \end{column}
    \begin{column}{0.5\textwidth}
      \textbf{Inconvénients}
      \begin{itemize}
        \item Dépendance externe pour une fonctionnalité que \texttt{fetch} couvre presque
        \item Poids supplémentaire dans le bundle
      \end{itemize}
    \end{column}
  \end{columns}
\end{frame}

% ================================
\section{API de données : Open Trivia DB}

\begin{frame}[fragile]{Open Trivia DB}
  \textbf{Qu'est-ce que c'est ?} API REST publique et gratuite fournissant des questions de quiz par catégorie/difficulté.

  \vspace{0.3cm}
  \textbf{Pourquoi ce choix ?}
  \begin{itemize}
    \item Gratuite, sans clé API, simple d'usage
    \item Permet de filtrer par \texttt{category}, \texttt{difficulty}, \texttt{type}
    \item Liste des catégories disponible via \texttt{api\_category.php}
  \end{itemize}

  \vspace{0.3cm}
  \textbf{Exemple d'appel} :
  \begin{lstlisting}
GET https://opentdb.com/api.php?amount=10&type=multiple
    &category=9&difficulty=easy
  \end{lstlisting}
\end{frame}

\begin{frame}{Open Trivia DB : avantages / inconvénients}
  \begin{columns}
    \begin{column}{0.5\textwidth}
      \textbf{Avantages}
      \begin{itemize}
        \item Aucune authentification requise
        \item Données prêtes à l'emploi (questions, réponses, mauvaises réponses)
        \item Filtrage simple par paramètres d'URL
      \end{itemize}
    \end{column}
    \begin{column}{0.5\textwidth}
      \textbf{Inconvénients}
      \begin{itemize}
        \item \textbf{Rate limiting} agressif (HTTP 429)
        \item Texte encodé en entités HTML (\texttt{\&quot;}, \texttt{\&amp;}...)
        \item Contenu uniquement en \textbf{anglais}
        \item Pas de garantie de disponibilité (service tiers)
      \end{itemize}
    \end{column}
  \end{columns}

  \vspace{0.3cm}
  \textbf{Ce que j'ai appris} : gérer les erreurs réseau (429), décoder le HTML reçu, et limiter le nombre d'appels (cache via \texttt{window.\_\_quizData}).
\end{frame}

% ================================
\section{État global : Context API}

\begin{frame}{Context API (React)}
  \textbf{Qu'est-ce que c'est ?} Mécanisme natif de React pour partager un état global sans \og prop drilling \fg{}.

  \vspace{0.3cm}
  \textbf{Pourquoi ce choix ?}
  \begin{itemize}
    \item Besoin de partager le \textbf{thème} et la \textbf{langue} dans toute l'app
    \item Pas besoin d'une librairie externe (Redux, Zustand...) pour 2 valeurs simples
    \item Couplé à \texttt{localStorage} pour persister les préférences
  \end{itemize}

  \vspace{0.3cm}
  \begin{columns}
    \begin{column}{0.5\textwidth}
      \textbf{Avantages}
      \begin{itemize}
        \item Intégré à React, zéro dépendance
        \item API simple : \texttt{Provider} + \texttt{useContext}
      \end{itemize}
    \end{column}
    \begin{column}{0.5\textwidth}
      \textbf{Inconvénients}
      \begin{itemize}
        \item Re-render de tous les consommateurs au changement
        \item Devient difficile à maintenir si l'état global grossit
      \end{itemize}
    \end{column}
  \end{columns}
\end{frame}

% ================================
\section{Theming \& i18n}

\begin{frame}[fragile]{Variables CSS pour le mode sombre}
  \textbf{Approche} : variables CSS (\texttt{--bg-color}, \texttt{--text-color}, \texttt{--accent-color}...) redéfinies sous \texttt{[data-theme="dark"]}.

  \begin{lstlisting}
:root { --bg-color: #f0f4ff; --text-color: #2c3e50; }
[data-theme="dark"] { --bg-color: #1a1d24; --text-color: #f0f0f0; }
  \end{lstlisting}

  \vspace{0.3cm}
  \begin{columns}
    \begin{column}{0.5\textwidth}
      \textbf{Avantages}
      \begin{itemize}
        \item Pas de librairie CSS-in-JS
        \item Changement de thème instantané (1 attribut HTML)
      \end{itemize}
    \end{column}
    \begin{column}{0.5\textwidth}
      \textbf{Inconvénients}
      \begin{itemize}
        \item Nécessite de revoir toutes les couleurs codées en dur
        \item Pas de typage / autocomplétion
      \end{itemize}
    \end{column}
  \end{columns}
\end{frame}

\begin{frame}{Internationalisation maison}
  \textbf{Approche} : dictionnaire JS (\texttt{translations.js}) + fonction \texttt{t(cle)} via \texttt{LanguageContext}.

  \vspace{0.3cm}
  \begin{columns}
    \begin{column}{0.5\textwidth}
      \textbf{Avantages}
      \begin{itemize}
        \item Très léger, aucune dépendance (vs. i18next)
        \item Suffisant pour 2 langues et peu de textes
      \end{itemize}
    \end{column}
    \begin{column}{0.5\textwidth}
      \textbf{Inconvénients}
      \begin{itemize}
        \item Pas de pluriels, formats de dates, etc.
        \item Le contenu de l'API (questions) reste en anglais
      \end{itemize}
    \end{column}
  \end{columns}

  \vspace{0.3cm}
  \textbf{Ce que j'ai appris} : structurer des clés imbriquées (\texttt{quiz.scorePerfect}) et garder la traduction synchronisée avec l'UI.
\end{frame}

% ================================
\section{Manette : Gamepad API}

\begin{frame}{Gamepad API}
  \textbf{Qu'est-ce que c'est ?} API web native (\texttt{navigator.getGamepads()}) pour lire l'état des manettes connectées.

  \vspace{0.3cm}
  \textbf{Pourquoi ce choix ?}
  \begin{itemize}
    \item Aucune librairie nécessaire, support natif des navigateurs
    \item Permet de naviguer/sélectionner les réponses sans clavier/souris
  \end{itemize}

  \vspace{0.3cm}
  \begin{columns}
    \begin{column}{0.5\textwidth}
      \textbf{Avantages}
      \begin{itemize}
        \item Expérience originale et ludique
        \item API simple, basée sur du \textit{polling}
      \end{itemize}
    \end{column}
    \begin{column}{0.5\textwidth}
      \textbf{Inconvénients}
      \begin{itemize}
        \item Pas d'événements : il faut \textit{poller} via \texttt{requestAnimationFrame}
        \item Mappings de boutons variables selon les manettes
        \item Gestion du \og cooldown \fg{} pour éviter les répétitions
      \end{itemize}
    \end{column}
  \end{columns}
\end{frame}

% ================================
\section{Bilan}

\begin{frame}{Ce que j'ai appris}
  \begin{itemize}
    \item Structurer une app React en pages / composants / contextes
    \item Gérer un état complexe (score, timer, historique) avec des hooks
    \item Déboguer des effets de bord et des problèmes de \textit{stale state}
    \item Consommer une API REST tierce avec ses contraintes (rate limiting, encodage)
    \item Mettre en place un système de thème et d'internationalisation léger, sans librairie lourde
    \item Utiliser une API navigateur peu courante (Gamepad)
  \end{itemize}
\end{frame}

\begin{frame}{Pistes d'évolution}
  \begin{itemize}
    \item Authentification (Firebase Auth)
    \item Sauvegarde des scores / classement (Firestore)
    \item Mode multijoueur, sons, animations
  \end{itemize}

  \vspace{1cm}
  \centering \Large Merci pour votre attention !
\end{frame}

\end{document}
